\section{Top-level Overview}

The proposed LoRa-based V2V mesh system uses a \textbf{three-layer architecture} so that application
features, network behavior, and radio transmission are kept separate. Keeping the system layered like
this makes it easier to build, test, and modify, while still supporting distributed safety-message
dissemination between vehicles over a constrained long-range wireless link.

The \textbf{Application Layer} is where vehicle services create and use messages. It generates outbound
traffic for neighboring vehicles and processes inbound messages delivered by the lower layers. We keep
this layer independent of routing policies and radio details, so service-level logic stays clean.

The \textbf{Network Layer} is the control core of the system. It manages delivery and forwarding
decisions by adding network metadata (such as hop counts, timestamps, and originator identifiers),
deciding whether a received packet should be delivered locally and/or forwarded, suppressing duplicates
to save channel usage, and scheduling relay timing to reduce contention and collisions. Keeping these
policies in one place also makes it easier to test and tune forwarding behavior.

The \textbf{Link Layer} sits at the radio boundary and handles actual transmission and reception on the
LoRa hardware. It packetizes serialized messages, manages radio state changes and channel sensing, and can
optionally use selective acknowledgment and retry for reliability. Since hardware details are contained here,
the upper layers do not need to change if the physical backend changes.

In normal operation, application payloads go down the stack to be transmitted, and received radio frames go
up through the Link and Network layers before reaching the Application Layer. To support both real hardware
runs and simulation-based testing, the Network Layer is kept hardware-independent using a clear abstract
interface to the Link Layer. This lets the same logic run either with real radio devices or with a
software-based channel simulator for controlled validation. The overall layering and data flow are shown in
Fig.~\ref{fig:system_architecture}.

\begin{figure}[h]
    \centering
    \includegraphics[width=0.45\textwidth]{"assets/architecture/top_level.jpg"}
    \caption{Three-layer architecture of the LoRa-based V2V mesh system, showing the Application, Network,
    and Link layers, the message flow for transmission and reception, peer vehicles, and the wireless channel.}
    \label{fig:system_architecture}
\end{figure}
